\documentclass[10pt, a4paper]{article}

\usepackage{amsmath}
\usepackage[a4paper, margin=1in]{geometry}
\usepackage{graphicx}
\usepackage{hyperref}
\usepackage{parskip}
\usepackage{svg}

\newcommand{\codeWord}[1]{\texttt{#1}}


\title{COMP0233: Assessment Framework}
\author{William Graham}
\date{\today}

\def\moduleGrade{g_{\mathrm{module}}}
\def\courseworkGrade{g_{\mathrm{cw}}}
\def\examGrade{g_{\mathrm{exam}}}

\begin{document}

\maketitle

\tableofcontents

% Redefine internal link colour, so the TOContents has normal text, but thereafter internal links are highlighted.
\hypersetup{
  linkcolor = magenta,
}

\section{TLDR: How your grade is calculated} \label{sec:tldr}

COMP0233 comprises two assessment components;

\begin{itemize}
  \item a group coursework component,
  \item an exam component,
\end{itemize}

whose individual grades are averaged to provide your overall grade for the module.
That is to say; if you achieve a grade of $\courseworkGrade \%$ in the coursework component, and $\examGrade \%$ in the exam, then your overall percentage mark for COMP0233 ($\moduleGrade$) will be

\begin{align*}
  \moduleGrade = \frac{\courseworkGrade + \examGrade}{2}.
\end{align*}

\section{UCL Assessment Framework}

In this section we will summarise the relevant portions of the UCL academic manual that pertain to assessments and the award of grades.
You can view the full academic manual online by \href{https://www.ucl.ac.uk/study/current-students/academic-manual}{heading to its website} - it should be your first point of call for information about your studies at UCL.

\subsection{Assessment Components and Submissions}

Modules at UCL are assessed through ``assessed components'', or just ``components''.
A module always has at least one component - most typically have either one or two, with COMP0233 being the latter (two).

A \emph{component} will require you to make a number of ``\emph{submissions}'' in order to be awarded a grade for that component.
Each submission should be understood as a standalone piece of the component; the most commonly seen example is the requirement to write a report (one submission) and then give a presentation on the subject of your report (another submission).
If any submissions are missing, you are illegible for a grade in its parent component - \emph{regardless} of whether all other submissions are provided.
Note that the term ``submission'' is often conflated with ``attempt'' - please be careful not to confuse these terms.
You are typically allowed multiple attempts at a submission (for example, a draft report before committing to submitting the actual report), and many support uploading iterative versions as you progress with the assessment.

You will be awarded a grade for each of the individual \emph{components} of a module.
These grades will then be combined via a weighted average into your overall module grade - for COMP0233, \hyperref[sec:tldr]{this is simply the average of the two components that make up the module}.
This means that failure (or non-submission) of a component does not necessarily mean failure of the module as a whole, but it will obviously hurt your overall chances.

\subsection{Pass Mark and Re-assessment}

COMP0233 follows central UCL policy concerning award of grades and accompanying CATS credit, re-assessment, and condonement.
Your first point of call for information about these policies should be the academic manual, or your personal tutor who can direct you to your department's assessment team.

Broadly speaking, you must achieve a module grade $\moduleGrade$ in COMP0233 of:

\begin{itemize}
  - 50\% to pass, if you are on a postgraduate degree.
  - 40\% to pass, if you are on an undergraduate degree.
  - If you fall short of the pass mark for your program, your department will consider whether you are eligible for condonement.
  If you are awarded a pass with condonement, treat yourself as if you had passed when reading on.
  The award of condonement is at the discretion of your department, so \emph{please do not ask questions about award of condonement to the module leads} - we cannot help you in this regard.
\end{itemize}

Assuming you have passed, your grade in COMP0233 will then contribute to your overall degree average in the manner dictated by your program.
Nothing else in this section is relevant in the event you pass.

If you fail the module, you may be required to take re-assessment in the Late Summer Assessment (LSA) period.
Here, there are a further three cases depending on the grade you achieved in each component of the module; your coursework grade $\courseworkGrade$ and your grade in the exam $\examGrade$:

\begin{itemize}
  \item f
\end{itemize}


\section{COMP0233: Module Assessment}

\subsection{Exam Component}

\subsection{Group Coursework Component}

% Figures \ref{fig:simpsons-interview} and \ref{fig:beaver-refactoring} were made using \href{https://imgflip.com/memegenerator}{ImgFlip}; \url{https://imgflip.com/memegenerator}.

\end{document}